\chapter{绪论}

\section{研究背景}

无线电自动调制识别（Automatic Modulation Classification, AMC）是在未知调制方式条件下，根据接收信号自动判别调制类别的任务。它是频谱感知、认知无线电、通信监测和非合作通信分析中的基础环节。RadioML 系列数据集和基于卷积神经网络的端到端识别方法推动了深度学习在 AMC 任务中的应用\cite{oshea2016radioml,oshea2016cnn,west2017deep}。与人工统计特征方法相比，深度模型能够直接从原始 I/Q 序列或派生时频表示中学习判别特征，但其结果也更依赖数据划分、训练随机性和评价协议。

传统 AMC 研究长期关注似然方法、高阶累积量和人工特征分类器\cite{azzouz1996amr,swami2000cumulants,dobre2007survey,hameed2009likelihood}，深度学习方法则把特征提取和分类器学习合并到端到端训练流程中。低信噪比条件是 AMC 任务中的主要困难之一。噪声会削弱 I/Q 序列中的相位、幅度和瞬时频率线索，使不同调制类别在观测空间中更加接近。短时傅里叶变换（Short-Time Fourier Transform, STFT）将信号局部频率结构显式展开\cite{oppenheim2010dtsp,allen1977shorttime,griffin1984stft,cohen1995timefrequency,boashash2015timefrequency}，理论上可能为低信噪比样本提供与原始 I/Q 序列互补的判别信息。

然而，现有深度 AMC 研究在实验协议上存在若干不一致：数据划分方式和训练随机性因文献而异，使不同工作的结论难以直接对比；多视图融合方法通常在未控制参数规模的条件下报告性能改善，难以区分视图融合收益与参数增量；低信噪比分组性能的改善常缺乏配对统计检验支撑。这些局限表明，需要在固定数据划分、统一模型矩阵和配对统计检验下对多视图融合方法进行系统评估。


\section{研究问题}

本文研究 \datasetname{} 上的 11 类调制分类问题。给定长度为 128 的双通道 I/Q 样本，模型输出对应的调制类别。任务范围聚焦于调制分类性能，信噪比估计、信道识别、频谱占用检测和端到端部署系统优化不纳入本报告评估。

围绕这一任务，本文关注三个问题：第一，在同一固定划分和相同训练种子集合下，I/Q 序列模型、STFT 时频模型和 I/Q--STFT 融合模型的整体表现如何；第二，STFT 视图是否在低信噪比样本上提供统计显著的局部补充信息；第三，门控融合和参数匹配控制是否能为融合收益提供可信的消融解释。

\section{本文工作与贡献}

本文在 \datasetname{} 上建立了一套协议化的多视图融合 AMC 实验框架，并在此基础上提出提案模型 \path{fusion_cldnn_stft}。实验矩阵包含 \path{cnn1d}、\path{resnet1d}、\path{tfcnn_stft}、\path{fusion_iq_stft}、\path{cldnn}、\path{mcldnn}、\path{lwamcnet}、\path{iq_param_matched} 和 \path{gated_fusion_iq_stft} 共 9 个模型行，每个模型使用 3 个训练种子，因此主实验包含 27 个运行单元。模型类型覆盖原始 I/Q 卷积基线、STFT 时频分支、静态多视图融合、门控多视图融合、卷积-循环混合模型、轻量模型和参数匹配 I/Q 控制模型。

本文的主要贡献如下：
\begin{enumerate}
  \item 建立了固定划分、三训练种子和 9 个模型行的协议化实验框架，通过参数匹配 I/Q 控制模型区分视图融合收益与参数规模增量，并完整保留门控融合和 MCLDNN 异常种子等负结果，确保比较结论不依赖结果筛选。
  \item 使用同一测试样本上的配对预测结果进行自助法置信区间估计和 McNemar 检验\cite{efron1979bootstrap,mcnemar1947note}，区分统计显著的发现与点估计差异，避免仅凭均值做外推结论。
  \item 通过四项独立扩展实验（延长训练预算、低信噪比加权交叉熵、扩展复杂度与单线程 CPU 延迟测量、低信噪比类别混淆后处理）依次排除主实验结果的常见替代解释，确认主要结论不是欠训练或损失函数设计不当造成的。
  \item 在文献校准基础上提出提案模型 \path{fusion_cldnn_stft}：以 CLDNN 风格 I/Q 编码器替换静态融合中的轻量 I/Q 分支，并结合相位旋转加循环时移信号域增强与标签平滑 0.1。在同一固定划分和三训练种子下，提案模型整体准确率超越主实验 \path{cldnn} 基线 1.35 个百分点，并在整体、低、中、高信噪比四项指标上同时正向。
\end{enumerate}

本文实验范围限于 \datasetname{} 数据集、固定划分 \splitname{} 和 3 个训练种子 \mainseeds{}；其他数据集、真实外场信道与端到端部署效率的泛化性留作后续工作。

\begin{figure}[H]
  \centering
  \resizebox{\textwidth}{!}{%
\begin{tikzpicture}[
  node distance=1.05cm and 1.1cm,
  stage/.style={rounded corners, draw=#1!75!black, fill=#1!10, very thick,
    align=center, minimum width=3.0cm, minimum height=1.05cm, font=\small},
  evidence/.style={rounded corners, draw=black!55, fill=black!4, thick,
    align=center, minimum width=3.0cm, minimum height=0.85cm, font=\small},
  arrow/.style={-{Latex[length=2.5mm]}, thick, draw=black!70}
]
\node[stage=blue] (data) {\datasetname{}\\220,000 个样本};
\node[stage=teal, right=of data] (split) {固定划分\\\splitname{}};
\node[stage=orange, right=of split] (models) {9 个模型行\\3 个训练种子};
\node[evidence, below=of models] (pred) {测试集预测归档\\27 个预测文件};
\node[evidence, left=of pred] (agg) {聚合结果表\\主表/低信噪比/延迟};
\node[evidence, right=of pred] (tests) {配对统计检验\\自助法 + McNemar};
\node[stage=green, below=of pred] (report) {课程报告结论\\受证据边界约束};
\draw[arrow] (data) -- (split);
\draw[arrow] (split) -- (models);
\draw[arrow] (models) -- (pred);
\draw[arrow] (pred) -- (tests);
\draw[arrow] (pred) -- (agg);
\draw[arrow] (agg) -- (report);
\draw[arrow] (tests) -- (report);
\node[evidence, below=0.55cm of agg, fill=gray!8] (exclude) {诊断性/工程检查材料\\不进入主结果};
\draw[arrow, dashed, draw=gray!70] (exclude) -- (report);
\end{tikzpicture}%
}

  \caption{实验证据链示意图。}
  \label{fig:protocol-pipeline}
\end{figure}


\section{全文结构}

第~2~章描述数据集、实验协议与评价方法。第~3~章介绍各模型结构、输入视图设计与提案模型。第~4~章报告主实验结果、扩展实验与统计检验。第~5~章讨论结论的解释范围与局限性。第~6~章总结全文并指出后续工作方向。
